\documentclass[12pt]{article}
\usepackage{ceai}

\begin{document}
\title{Tutorial 04:  ML Production in Action - 02}
\author{}
\date{}

\maketitle

\section{Part B: Nonlinear Classification Models}

In Part A, we formulated a binary classification problem using the UCI concrete dataset and evaluated linear classifiers such as logistic regression and perceptron models.  

In this part, we extend the study to more expressive machine learning models capable of learning nonlinear decision boundaries.

The objectives of this tutorial are:

\begin{itemize}
\item Train a Multi-Layer Perceptron (MLP) classifier
\item Understand regularization in neural networks
\item Train Support Vector Machine (SVM) classifiers
\item Study the role of kernel functions
\item Evaluate classifier capability using ROC and Precision–Recall curves
\item Introduce the concept of model selection and hyperparameters
\end{itemize}

All experiments will be implemented using Python in a Google Colab environment with the \texttt{scikit-learn} library.

\section{Dataset and Problem Definition}

We continue using the UCI Concrete Compressive Strength dataset introduced in Part A.

The classification labels are defined based on the estimated 28-day compressive strength:

\[
y =
\begin{cases}
1 & f'_c(28) \le 40 \text{ MPa (NSC)} \\
0 & f'_c(28) > 40 \text{ MPa (HSC)}
\end{cases}
\]

The dataset contains mixture composition features such as cement, water, slag, fly ash, and age.

\subsection{Data Preparation}

\begin{py-box}
import pandas as pd
import numpy as np

from sklearn.model_selection import train_test_split
from sklearn.preprocessing import StandardScaler

df = pd.read_excel("Concrete_Data.xlsx")

X = df.iloc[:,0:8]
strength = df.iloc[:,8]

y = (strength <= 40).astype(int)

X_train, X_test, y_train, y_test = train_test_split(
    X, y, test_size=0.25, random_state=42)

scaler = StandardScaler()

X_train = scaler.fit_transform(X_train)
X_test = scaler.transform(X_test)
\end{py-box}

Feature standardization is important for both neural networks and SVM models because these algorithms rely on geometric distances in feature space.

\section{Multi-Layer Perceptron (MLP)}

\subsection{Model Structure}

A Multi-Layer Perceptron is a feedforward neural network composed of stacked layers of neurons.

A simple one-hidden-layer network can be written as

\[
h = \sigma(W_1 x + b_1)
\]

\[
\hat{y} = \sigma(W_2 h + b_2)
\]

where

\begin{itemize}
\item $x$ is the feature vector
\item $W_1, W_2$ are weight matrices
\item $\sigma(\cdot)$ is a nonlinear activation function
\end{itemize}

Common activation functions include:

\begin{itemize}
\item ReLU
\item Sigmoid
\item Hyperbolic tangent
\end{itemize}

\subsection{Training an MLP Classifier}

\begin{py-box}
from sklearn.neural_network import MLPClassifier
from sklearn.metrics import accuracy_score

mlp = MLPClassifier(
    hidden_layer_sizes=(20,10),
    activation='relu',
    max_iter=2000,
    random_state=0
)

mlp.fit(X_train, y_train)

y_pred = mlp.predict(X_test)

print("Test accuracy:", accuracy_score(y_test, y_pred))
\end{py-box}

\subsection{Regularization in Neural Networks}

Neural networks often contain many parameters and can easily overfit the training data.

Regularization introduces a penalty term to the optimization objective:

\[
L = L_{data} + \lambda ||W||^2
\]

where

\begin{itemize}
\item $L_{data}$ is the classification loss
\item $||W||^2$ penalizes large weights
\item $\lambda$ controls regularization strength
\end{itemize}

In \texttt{scikit-learn}, this parameter is called \texttt{alpha}.

\begin{py-box}
mlp_reg = MLPClassifier(
    hidden_layer_sizes=(20,10),
    activation='relu',
    alpha=0.01,
    max_iter=2000
)

mlp_reg.fit(X_train, y_train)

print("Regularized accuracy:", mlp_reg.score(X_test, y_test))
\end{py-box}

\subsection{Decision Boundary of an MLP}

The trained network defines a decision function

\[
g(x) = f(u(x)|w)
\]

where $w$ denotes all trained weights.

For binary classification, the last layer often uses the sigmoid function

\[
\sigma(z)=\frac{1}{1+e^{-z}}
\]

which outputs a probability estimate.

The decision boundary corresponds to

\[
\sigma(z)=0.5
\]

or equivalently

\[
z=0
\]

Because the dataset has multiple dimensions, we visualize the decision boundary using two selected features.

\subsection{Training a 2D Visualization Model}

\begin{py-box}
import matplotlib.pyplot as plt

X2 = X[['Cement (component 1)(kg in a m^3 mixture)',
        'Water  (component 4)(kg in a m^3 mixture)']]

y = (strength <= 40).astype(int)

scaler2 = StandardScaler()
X2_scaled = scaler2.fit_transform(X2)

mlp2 = MLPClassifier(
    hidden_layer_sizes=(10,10),
    activation='relu',
    max_iter=2000
)

mlp2.fit(X2_scaled, y)
\end{py-box}

\subsection{Plotting the Decision Boundary}

\begin{py-box}
x_min, x_max = X2_scaled[:,0].min()-1, X2_scaled[:,0].max()+1
y_min, y_max = X2_scaled[:,1].min()-1, X2_scaled[:,1].max()+1

xx, yy = np.meshgrid(
    np.linspace(x_min,x_max,200),
    np.linspace(y_min,y_max,200)
)

grid = np.c_[xx.ravel(),yy.ravel()]

prob = mlp2.predict_proba(grid)[:,1]

Z = prob.reshape(xx.shape)

plt.contourf(xx,yy,Z,levels=20,cmap='coolwarm',alpha=0.6)

plt.scatter(X2_scaled[:,0],X2_scaled[:,1],c=y,cmap='coolwarm')

plt.title("MLP Decision Boundary")
plt.show()
\end{py-box}

\section{Comparing Decision Boundaries}

To understand the capability of different models, we compare the decision boundaries of three classifiers:

\begin{itemize}
\item Logistic regression
\item Multi-layer perceptron
\item Kernel SVM
\end{itemize}

\begin{py-box}
from sklearn.linear_model import LogisticRegression
from sklearn.svm import SVC

log_model = LogisticRegression()
log_model.fit(X2_scaled,y)

svm_rbf = SVC(kernel='rbf',probability=True)
svm_rbf.fit(X2_scaled,y)
\end{py-box}

\begin{py-box}
models = [log_model, mlp2, svm_rbf]

titles = [
"Logistic Regression",
"MLP",
"Kernel SVM"
]

plt.figure(figsize=(12,4))

for i,model in enumerate(models):

    plt.subplot(1,3,i+1)

    prob = model.predict_proba(grid)[:,1]
    Z = prob.reshape(xx.shape)

    plt.contourf(xx,yy,Z,levels=20,cmap='coolwarm',alpha=0.6)
    plt.scatter(X2_scaled[:,0],X2_scaled[:,1],c=y,cmap='coolwarm')

    plt.title(titles[i])

plt.show()
\end{py-box}

This comparison illustrates how model complexity affects the geometry of the decision boundary.

\section{Support Vector Machines}

\subsection{Linear Support Vector Machine}

A Support Vector Machine (SVM) seeks a separating hyperplane that maximizes the geometric margin between two classes.

A linear decision function has the form

\[
f(x) = w^T x + b
\]

where

\begin{itemize}
\item $w$ is the weight vector
\item $b$ is the bias
\end{itemize}

The classification rule is

\[
\hat{y} =
\begin{cases}
1 & f(x) \ge 0 \\
0 & f(x) < 0
\end{cases}
\]

The decision boundary corresponds to

\[
w^T x + b = 0
\]

\paragraph{Margin Definition}

The margin is the distance between the decision boundary and the closest training samples.

For a linear classifier, the margin width is

\[
\text{margin width} = \frac{2}{||w||}
\]

Thus maximizing the margin is equivalent to minimizing the norm of the weight vector.

\paragraph{Soft-Margin Optimization}

In practice, perfect separation is rarely possible.  
The \emph{soft-margin SVM} introduces slack variables $\xi_i$ that allow some misclassification.

The optimization problem becomes

\[
\min_{w,b,\xi}
\quad
\frac{1}{2}||w||^2
+
C\sum_{i=1}^{n}\xi_i
\]

subject to

\[
y_i(w^T x_i + b) \ge 1 - \xi_i
\]

\[
\xi_i \ge 0
\]

The parameter $C$ controls the trade-off between

\begin{itemize}
\item maximizing the margin (small $||w||$)
\item penalizing classification errors
\end{itemize}

\paragraph{Interpretation of $C$}

\begin{itemize}
\item Small $C$  
      \begin{itemize}
      \item larger margin
      \item more tolerance for classification errors
      \item stronger regularization
      \end{itemize}

\item Large $C$  
      \begin{itemize}
      \item narrower margin
      \item fewer training errors
      \item weaker regularization
      \end{itemize}
\end{itemize}

\subsection*{Training a Linear SVM}

\begin{py-box}
from sklearn.svm import SVC

svm_linear = SVC(kernel='linear', C=1)

svm_linear.fit(X_train, y_train)

print("Linear SVM accuracy:",
      svm_linear.score(X_test, y_test))
\end{py-box}


\subsection{Effect of the Parameter C}

The parameter $C$ controls the trade-off between margin width and training errors.

\begin{py-box}
for C in [0.01,0.1,1,10,100]:

    svm = SVC(kernel='linear',C=C)

    svm.fit(X_train,y_train)

    acc = svm.score(X_test,y_test)

    print("C =",C,"accuracy =",acc)
\end{py-box}

\section{Kernel Functions}

\subsection{Polynomial Kernel}

\[
K(x_i,x_j)=(x_i^Tx_j+c)^d
\]

\begin{py-box}
svm_poly = SVC(kernel='poly',degree=2)

svm_poly.fit(X_train,y_train)

print("Polynomial kernel accuracy:",
svm_poly.score(X_test,y_test))
\end{py-box}

\subsection{Radial Basis Function Kernel}

\[
K(x_i,x_j)=
\exp(-\gamma||x_i-x_j||^2)
\]

\begin{py-box}
svm_rbf = SVC(kernel='rbf')

svm_rbf.fit(X_train,y_train)

print("RBF kernel accuracy:",
svm_rbf.score(X_test,y_test))
\end{py-box}


\subsection{Remark: Relation Between MLP Regularization and the SVM Margin}

Earlier in this tutorial we introduced regularization for neural networks.

The regularized training objective for a neural network can be written as

\[
L = L_{\text{data}} + \lambda ||W||^2
\]

where

\begin{itemize}
\item $L_{\text{data}}$ is the classification loss
\item $||W||^2$ penalizes large weights
\item $\lambda$ controls the strength of regularization
\end{itemize}

This formulation is conceptually similar to the SVM objective:

\[
\min
\quad
\frac12 ||w||^2
+
C \sum \xi_i
\]

Both formulations contain two competing objectives:

\begin{itemize}
\item fitting the training data
\item controlling model complexity
\end{itemize}

However, the mechanisms differ.

\paragraph{Neural Network Regularization}

\begin{itemize}
\item penalizes large network weights
\item encourages smoother decision surfaces
\item reduces overfitting
\end{itemize}

\paragraph{SVM Regularization}

\begin{itemize}
\item explicitly maximizes the margin between classes
\item the parameter $C$ controls how strictly training points must satisfy the margin constraint
\end{itemize}

Thus we can interpret

\[
\lambda \quad \leftrightarrow \quad \frac{1}{C}
\]

as parameters that regulate the trade-off between

\begin{itemize}
\item data fitting
\item model complexity
\end{itemize}

In geometric terms:

\begin{itemize}
\item neural network regularization restricts the magnitude of the weights that shape the decision surface
\item SVM directly controls the width of the separating margin
\end{itemize}

Both approaches therefore serve a similar goal:  
\emph{improving generalization by preventing overly complex decision boundaries.}
\section{Evaluating Classifier Capability}

Accuracy alone does not fully describe classifier performance.

Two important evaluation tools are:

\begin{itemize}
\item ROC curve
\item Precision–Recall curve
\end{itemize}



\subsection{ROC Curve}

The ROC curve plots

\begin{itemize}
\item True Positive Rate
\item False Positive Rate
\end{itemize}

across classification thresholds.

\subsection{Precision–Recall Curve}

The PR curve plots

\begin{itemize}
\item Precision
\item Recall
\end{itemize}

which is particularly informative for imbalanced datasets.

\subsection{Generating ROC and PR Curves}

\begin{py-box}
from sklearn.metrics import roc_curve
from sklearn.metrics import precision_recall_curve

scores = svm_rbf.decision_function(X_test)

fpr,tpr,_ = roc_curve(y_test,scores)

precision,recall,_ = precision_recall_curve(y_test,scores)

plt.plot(fpr,tpr)
plt.title("ROC Curve")
plt.show()

plt.plot(recall,precision)
plt.title("Precision Recall Curve")
plt.show()
\end{py-box}

\section{Model Selection}

Throughout this tutorial we used default hyperparameters:

\begin{itemize}
\item number of hidden layers
\item number of neurons
\item activation functions
\item SVM parameter $C$
\item kernel type
\end{itemize}

Selecting these values appropriately is called the \textbf{model selection problem}.

Model selection may involve:

\begin{itemize}
\item tuning hyperparameters within a model
\item choosing between different model families
\end{itemize}

Typical techniques include:

\begin{itemize}
\item cross validation
\item grid search
\item automated hyperparameter optimization
\end{itemize}

These methods will be explored in later tutorials.
\end{document}
